\begin{statementbox}
	\textbf{\textcolor{CStatement}{条件}}
	\begin{description}[style=nextline, leftmargin=2.8em, labelsep=0.5em, itemsep=0.35em]
		\item[\textbf{\textcolor{CStatement}{条件 1（三角形与中点）：}}] 在$\triangle ABC$中，$BC=a$, $CA=b$, $AB=c$，$D,E,F$分别为$BC,CA,AB$的中点，中线$AD=m_a$，重心为$G$.
	\end{description}

	\textbf{\textcolor{CStatement}{结论}}
	\begin{description}[style=nextline, leftmargin=2.8em, labelsep=0.5em, itemsep=0.45em]
		\item[\textbf{\textcolor{CStatement}{结论一（中线长公式）：}}]
			\textbf{\textcolor{CStatement}{直观描述：}}
			已知三边可求中线长度.
			\[
				m_a^2 = \frac12(b^2+c^2) - \frac14 a^2.
			\]

		\item[\textbf{\textcolor{CStatement}{结论二（中位线定理）：}}]
			\textbf{\textcolor{CStatement}{直观描述：}}
			中点连线平行于底边且为半长.
			\[
				EF \parallel BC,\quad EF = \frac12 a.
			\]

		\item[\textbf{\textcolor{CStatement}{结论三（斜边中线定理）：}}]
			\textbf{\textcolor{CStatement}{直观描述：}}
			直角三角形斜边中线等于斜边一半.
			\[
				\angle A = 90^\circ \;\Rightarrow\; AD = \frac12 a.
			\]
	\end{description}

	\textbf{\textcolor{CStatement}{推广结论}}
	\begin{description}[style=nextline, leftmargin=2.8em, labelsep=0.5em, itemsep=0.45em]
		\item[\textbf{\textcolor{CStatement}{推广结论（重心向量性质）：}}]
			\textbf{\textcolor{CStatement}{直观描述：}}
			重心到三顶点向量和为零，分中线2:1.
			\[
				\overrightarrow{GA} + \overrightarrow{GB} + \overrightarrow{GC} = \mathbf{0},\quad AG : GD = 2 : 1.
			\]
	\end{description}

	\textbf{\textcolor{CStatement}{使用提醒（构造技巧）：}}
	倍长中线法：延长$AD$至$A'$使$AD=DA'$，则四边形$ABA'C$是平行四边形，且$A'B=AC$.
\end{statementbox}
